%======================================================================
% Chapter 12: Where to Go Next
%======================================================================
\chapter{Where to Go Next}\label{ch:next}

\begin{motivation}
You began this book with counting, fractions and the order of operations. You
have finished it with the mathematics of backpropagation, cross-entropy and
attention. That is a long road, and you now have the prerequisites that
textbooks and courses on machine learning take for granted. This short chapter
looks back at the big ideas, suggests projects that turn knowledge into skill,
sketches the mathematics that lies beyond these notes, and recommends books,
courses and videos for the next stage of the journey.
\end{motivation}

%----------------------------------------------------------------------
\section{The big ideas}\label{sec:big-ideas}
%----------------------------------------------------------------------

Beneath the many definitions and formulas, a handful of ideas carry most of the
weight. If you remember nothing else, remember these.

\begin{keytakeaways}[Seven big ideas]
\begin{enumerate}[label=\arabic*.]
  \item \textbf{Data are vectors, models are functions, layers are matrices.}
    Anything a computer learns from (an image, a sentence, a customer) becomes a
    list of numbers, and a model is a function of those numbers with adjustable
    parameters (\cref{ch:functions,ch:linalg1}).
  \item \textbf{Geometry reveals structure.} Lengths, angles and projections
    measure similarity and error; eigenvectors and the SVD find the most important
    directions in data (\cref{ch:geometry,ch:linalg2}).
  \item \textbf{Learning is optimisation.} Training means choosing the parameters
    that minimise a loss, usually by following the negative gradient downhill
    (\cref{ch:optimization}).
  \item \textbf{Derivatives measure sensitivity, and the chain rule composes them.}
    Backpropagation is the chain rule, organised so that all the derivatives of a
    huge network cost little more than evaluating it
    (\cref{ch:calculus,ch:multivariable,ch:dl}).
  \item \textbf{Uncertainty is described by probability.} Models output
    distributions, and the standard losses are negative log-likelihoods
    (\cref{ch:probability,ch:statistics}).
  \item \textbf{What matters is new data.} Fitting the training set is easy; the
    goal is generalisation, governed by the trade-off between bias and variance and
    measured honestly on held-out data (\cref{ch:statistics}).
  \item \textbf{Numbers on a computer need care.} Work with log-probabilities,
    subtract the maximum before exponentiating, scale initial weights to keep
    variances stable, and check gradients numerically
    (\cref{ch:functions,ch:calculus,ch:dl}).
\end{enumerate}
\end{keytakeaways}

%----------------------------------------------------------------------
\section{Consolidate by building}\label{sec:projects}
%----------------------------------------------------------------------

Mathematics becomes yours when you use it. The best next step is to implement
the methods of this book yourself, in Python with NumPy, checking every result
against the theory. The following projects are arranged roughly in order of
difficulty; each lists the chapters it relies on.

\begin{enumerate}[label=\textbf{Project \arabic*.},leftmargin=*,itemsep=5pt]
  \item \textbf{Linear regression three ways.} Fit a line to a real data set (for
    example, house prices against size) with the normal equations, with gradient
    descent and with stochastic gradient descent. Check that all three agree, and
    plot the loss against the number of steps for several learning rates.
    (\Cref{ch:linalg2,ch:optimization}.)
  \item \textbf{PCA on images.} Load a small data set of images of handwritten
    digits (one is included in the scikit-learn library), compute the principal
    components with the SVD, plot the first few as images, and reconstruct digits
    from $5$, $20$ and $50$ components. (\Cref{ch:linalg2}.)
  \item \textbf{Logistic regression with gradient checking.} Implement binary
    logistic regression with the gradient $(q - y)\vx$, verify the gradient with
    central differences, add $L^2$ regularisation, and report the test accuracy
    with a $95\%$ confidence interval. (\Cref{ch:calculus,ch:statistics,ch:dl}.)
  \item \textbf{Simulate the limit theorems.} Check the law of large numbers and
    the central limit theorem by simulation, for dice, for a skewed distribution
    such as the exponential, and for a distribution whose mean and variance are
    infinite (try $1/U$ for $U$ uniform on $[0, 1]$, and see what goes wrong). (\Cref{ch:probability,ch:statistics}.)
  \item \textbf{Grow the autodiff engine.} Extend the \texttt{Value} class of
    \cref{sec:autodiff} with subtraction, division, powers, $\exp$, $\ln$ and
    $\tanh$; build a small network from it and train it on the circle data of
    \cref{sec:training-loop}. (\Cref{ch:dl}.)
  \item \textbf{A digit classifier from scratch.} Train a network with one or two
    hidden layers on the MNIST digits, with softmax and cross-entropy, mini-batches
    and Adam, all in NumPy. Plot training and validation curves, experiment with
    initialisation and learning rates, and look at the digits it gets wrong.
    (\Cref{ch:optimization,ch:statistics,ch:dl}.)
  \item \textbf{Attention by hand.} Implement scaled dot-product attention for a
    short sequence of random vectors, verify that each row of weights sums to $1$,
    and observe what happens to the weights without the factor $1/\sqrt d$ as $d$
    grows. (\Cref{ch:probability,ch:dl}.)
\end{enumerate}

\begin{remark}[How to read mathematics in papers and textbooks]
When you meet an unfamiliar formula, do what this book has done throughout:
write down the shape of every vector and matrix; try the formula on a tiny
example with numbers ($2\times2$ matrices, three data points); draw a picture; and
ask what happens in extreme cases (a parameter equal to $0$, or very large).
\Cref{app:notation} lists the notation used here, which is standard in the
machine learning literature.
\end{remark}

%----------------------------------------------------------------------
\section{Mathematics beyond these notes}\label{sec:beyond}
%----------------------------------------------------------------------

These notes cover the prerequisites. The following topics build directly on
them and appear in more advanced work.

\begin{description}[style=nextline,leftmargin=1.5em,itemsep=4pt]
  \item[More linear algebra] Further matrix factorisations (LU, QR, Cholesky),
    numerical stability and the condition number, sparse matrices, and tensors
    (arrays with more than two indices), which are the basic objects of deep
    learning libraries. Builds on \cref{ch:linalg1,ch:linalg2}.
  \item[Fourier analysis and convolution] Representing signals as sums of sines
    and cosines; convolution, the operation behind convolutional neural networks
    for images and audio. Builds on \cref{ch:geometry,ch:calculus}.
  \item[Real analysis] Rigorous treatment of limits, sequences, series and
    continuity, needed to read proofs that algorithms converge. Builds on
    \cref{ch:calculus}.
  \item[Convex optimisation and duality] The Karush--Kuhn--Tucker conditions,
    which extend Lagrange multipliers to inequality constraints, duality, and
    convergence rates; support vector machines are the classic application.
    Builds on \cref{ch:optimization}.
  \item[Probability theory and stochastic processes] Measure-theoretic
    foundations; Markov chains and Markov chain Monte Carlo sampling; stochastic
    differential equations, which describe the noising process of diffusion
    models. Builds on \cref{ch:probability}.
  \item[Bayesian machine learning] Conjugate priors, Gaussian processes,
    variational inference and probabilistic graphical models. Builds on
    \cref{ch:probability,ch:statistics}.
  \item[Statistical learning theory] Why and when learning from a finite sample
    works: VC dimension, PAC learning and generalisation bounds. Builds on
    \cref{ch:statistics}.
  \item[Information theory] Mutual information, coding theorems and
    rate--distortion theory, which appear in representation learning and
    compression. Builds on \cref{sec:information-theory}.
  \item[Reinforcement learning] Markov decision processes, Bellman equations and
    policy gradients, for agents that learn from rewards. Builds on
    \cref{ch:probability,ch:optimization}.
  \item[Graphs] Graph theory and the spectral analysis of graphs, the basis of
    graph neural networks. Builds on \cref{ch:linalg2}.
\end{description}

%----------------------------------------------------------------------
\section{Recommended resources}\label{sec:reading}
%----------------------------------------------------------------------

There are excellent books and courses at every level, many of them free. The
following selection is organised by purpose; \cref{fig:roadmap} suggests one path
through it.

\begin{figure}[htbp]
\centering
\begin{tikzpicture}[>={Stealth[length=2.5mm]},
  stage/.style={draw,thick,rounded corners,align=left,font=\small,
    text width=11.2cm,inner sep=6pt},
  branch/.style={draw,thick,rounded corners,align=center,font=\small,
    text width=2.75cm,inner sep=4pt,minimum height=1.55cm}]
  \node[stage,fill=MLdef!10] (s1) at (0,0)
    {\textbf{1. Consolidate.} Redo the exercises; build the projects of
     \cref{sec:projects}; watch 3Blue1Brown's visual series on linear algebra
     and calculus.};
  \node[stage,fill=MLexample!10,below=0.6cm of s1] (s2)
    {\textbf{2. Machine learning foundations.} \emph{An Introduction to Statistical
     Learning}; \emph{Mathematics for Machine Learning}; an introductory machine
     learning course.};
  \node[stage,fill=MLintuition!14,below=0.6cm of s2] (s3)
    {\textbf{3. Deep learning.} \emph{Understanding Deep Learning};
     \emph{Dive into Deep Learning}; Karpathy's \emph{Zero to Hero}; fast.ai.};
  \draw[->,thick] (s1) -- (s2);
  \draw[->,thick] (s2) -- (s3);
  \node[branch,fill=MLml!10,below=1.0cm of s3.south west,anchor=north west] (b1)
    {\textbf{Probabilistic ML}\\Murphy; Bishop};
  \node[branch,fill=MLml!10,right=0.12cm of b1] (b2)
    {\textbf{Reinforcement learning}\\Sutton \& Barto};
  \node[branch,fill=MLml!10,right=0.12cm of b2] (b3)
    {\textbf{Theory}\\Shalev-Shwartz \& Ben-David};
  \node[branch,fill=MLml!10,right=0.12cm of b3] (b4)
    {\textbf{Research}\\reading and reproducing papers};
  \node[font=\small\bfseries,anchor=south west] at ([yshift=2pt]b1.north west) {4. Specialise};
  \foreach \b in {b1,b2,b3,b4}{\draw[->,thick] (s3.south) -- (\b.north);}
\end{tikzpicture}
\caption{A suggested path for further study.}
\label{fig:roadmap}
\end{figure}

\subsection{Mathematics}

\begin{itemize}[itemsep=3pt]
  \item Marc Peter Deisenroth, A.~Aldo Faisal and Cheng Soon Ong,
    \emph{Mathematics for Machine Learning} (Cambridge University Press, 2020;
    free PDF online). The natural next book: the same topics as
    \cref{ch:linalg1,ch:linalg2,ch:multivariable,ch:optimization,ch:probability}
    at a faster pace and greater depth, followed by applications to regression,
    PCA, Gaussian mixture models and support vector machines.
  \item Gilbert Strang, \emph{Introduction to Linear Algebra} (Wellesley-Cambridge
    Press), with his celebrated MIT 18.06 video lectures (free on MIT
    OpenCourseWare). Intuitive and applied. His \emph{Linear Algebra and Learning
    from Data} connects linear algebra to machine learning directly.
  \item Stephen Boyd and Lieven Vandenberghe, \emph{Introduction to Applied
    Linear Algebra: Vectors, Matrices, and Least Squares} (Cambridge University
    Press, 2018; free online). Gentle and full of applications.
  \item Sheldon Axler, \emph{Linear Algebra Done Right} (Springer; the fourth
    edition is open access). For readers who want a proof-based treatment.
  \item 3Blue1Brown (Grant Sanderson), \emph{Essence of Linear Algebra} and
    \emph{Essence of Calculus} (free video series). Unrivalled geometric
    intuition; ideal alongside \cref{ch:linalg1,ch:calculus}.
  \item Joseph K. Blitzstein and Jessica Hwang, \emph{Introduction to Probability}
    (CRC Press, 2nd edition 2019), with the free Harvard Stat 110 lectures. A
    superb, example-driven course in probability.
  \item Stephen Boyd and Lieven Vandenberghe, \emph{Convex Optimization}
    (Cambridge University Press, 2004; free PDF). The standard reference on
    convexity, duality and optimisation algorithms.
  \item David J.\,C. MacKay, \emph{Information Theory, Inference, and Learning
    Algorithms} (Cambridge University Press, 2003; free PDF). A delightful book
    linking information theory, Bayesian inference and neural networks.
\end{itemize}

\subsection{Machine learning}

\begin{itemize}[itemsep=3pt]
  \item Gareth James, Daniela Witten, Trevor Hastie and Robert Tibshirani,
    \emph{An Introduction to Statistical Learning} (Springer; editions with labs in
    R and in Python; free PDF). The friendliest introduction to the classical
    methods. Its big sibling, \emph{The Elements of Statistical Learning} by
    Hastie, Tibshirani and Friedman, is more mathematical.
  \item Christopher M. Bishop, \emph{Pattern Recognition and Machine Learning}
    (Springer, 2006). A classic, with a consistently probabilistic viewpoint.
  \item Kevin P. Murphy, \emph{Probabilistic Machine Learning: An Introduction}
    (MIT Press, 2022) and \emph{Advanced Topics} (2023); free drafts online.
    Comprehensive and up to date.
  \item Shai Shalev-Shwartz and Shai Ben-David, \emph{Understanding Machine
    Learning: From Theory to Algorithms} (Cambridge University Press, 2014; free
    online). The theory of generalisation.
  \item Richard S. Sutton and Andrew G. Barto, \emph{Reinforcement Learning: An
    Introduction} (MIT Press, 2nd edition 2018; free online).
\end{itemize}

\subsection{Deep learning}

\begin{itemize}[itemsep=3pt]
  \item Simon J.\,D. Prince, \emph{Understanding Deep Learning} (MIT Press, 2023;
    free PDF). Clear explanations and outstanding figures; the best next step
    after \cref{ch:dl}.
  \item Christopher M. Bishop and Hugh Bishop, \emph{Deep Learning: Foundations
    and Concepts} (Springer, 2024). A thorough modern textbook.
  \item Ian Goodfellow, Yoshua Bengio and Aaron Courville, \emph{Deep Learning}
    (MIT Press, 2016; free online). Its first part reviews linear algebra,
    probability, information theory and numerical computation.
  \item Aston Zhang, Zachary C. Lipton, Mu Li and Alexander J. Smola, \emph{Dive
    into Deep Learning} (free online at \texttt{d2l.ai}). Every idea is accompanied
    by runnable code.
  \item Andrej Karpathy, \emph{Neural Networks: Zero to Hero} (free video series).
    Starts by building an autodiff engine much like the one in \cref{sec:autodiff}
    and ends with a small GPT-style language model.
  \item fast.ai, \emph{Practical Deep Learning for Coders} (free course). A
    top-down, hands-on complement to the bottom-up approach of these notes.
\end{itemize}

%----------------------------------------------------------------------
\section{A final word}\label{sec:final-word}
%----------------------------------------------------------------------

Nobody reads mathematics once and understands it fully. Expect to return to
these chapters: when a paper writes $\nabla_{\vtheta}\E_{\vx\sim p}[\ell(\vx;
\vtheta)]$ and you need to recall why the expectation and the gradient can be
swapped for a mini-batch estimate; when a training run diverges and you suspect the
learning rate or the conditioning; when a model is confidently wrong and you want
to reason about its probabilities. Each time, the ideas will feel more familiar
and the formulas more like a language you speak.

Mistakes and confusion are part of the process, not signs of failure. Work small
examples, draw pictures, check shapes, and write code. Machine learning rewards
exactly the habits this book has tried to build: careful definitions, concrete
numbers before abstraction, and the patience to ask \emph{why}.

Good luck, and enjoy the journey.
